% GENERATED FILE. Edit canonical lessons, then run npm run generate:books.
% canonical-source-hash: fnv1a64:f1dd00d4592cd51f
% canonical-lessons: BN-C143-shoktishali, BN-C143-susbadu, BN-C143-borosoro, BN-C143-chhotokhato, BN-C143-hasikhushi

\chapter{Powerful, Tasty, Fairly big, Small, Cheerful}
\label{ch:bn-powerful-tasty-fairly-big-small-cheerful}
\hlchaptermodality{143}

\begin{chapteropening}
\textbf{By the end of this chapter:} \emph{I can say powerful, tasty, fairly big, small and cheerful.}

Complete the last lesson of chapter 143: i can say powerful, tasty, fairly big, small and cheerful.
\end{chapteropening}

\section[\texorpdfstring{\bn{শক্তিশালী} (shôktishāli)}{শক্তিশালী (shôktishāli)}]{\bn{শক্তিশালী} (shôktishāli) — powerful}

\label{lesson:BN-C143-shoktishali}



\begin{quote}
Before the new one: say the Bengali for sturdy, then the Bengali for weak.
\end{quote}

\subsection*{You'll want to know: \bn{শক্তিশালী}}
\textbf{\bn{শক্তিশালী}} — \emph{shôktishāli} — \textquotedblleft{}powerful\textquotedblright{}.

\begin{grammarlens}[title={What you've built}]
One more word for this chapter.
\end{grammarlens}

\subsection*{Your turn}
\begin{itemize}
\raggedright
  \item \emph{Say it:} \emph{shôktishāli}
  \item \emph{Say it:} \emph{shôktishāli}, once more
  \item \emph{Say it:} say \emph{shôktishāli}
  \item \emph{Recall:} say the Bengali for sturdy, then the Bengali for weak, then say \emph{shôktishāli} again
\end{itemize}

\begin{tcolorbox}[breakable,colback=teal!4,colframe=teal!35!black,title={Before you move on}]
What does \bn{শক্তিশালী} mean? (\textquotedblleft{}Powerful\textquotedblright{}.) Say it once more.
\end{tcolorbox}
\section[\texorpdfstring{\bn{সুস্বাদু} (susbādu)}{সুস্বাদু (susbādu)}]{\bn{সুস্বাদু} (susbādu) — tasty}

\label{lesson:BN-C143-susbadu}



\begin{quote}
Before the new one: say the Bengali for weak, then the Bengali for powerful.
\end{quote}

\subsection*{You'll want to know: \bn{সুস্বাদু}}
\textbf{\bn{সুস্বাদু}} — \emph{susbādu} — \textquotedblleft{}tasty\textquotedblright{}.

\begin{grammarlens}[title={What you've built}]
One more word for this chapter.
\end{grammarlens}

\subsection*{Your turn}
\begin{itemize}
\raggedright
  \item \emph{Say it:} \emph{susbādu}
  \item \emph{Say it:} \emph{susbādu}, once more
  \item \emph{Say it:} say \emph{susbādu}
  \item \emph{Recall:} say the Bengali for weak, then the Bengali for powerful, then say \emph{susbādu} again
\end{itemize}

\begin{tcolorbox}[breakable,colback=teal!4,colframe=teal!35!black,title={Before you move on}]
What does \bn{সুস্বাদু} mean? (\textquotedblleft{}Tasty\textquotedblright{}.) Say it once more.
\end{tcolorbox}
\section[\texorpdfstring{\bn{বড়সড়} (bôṛôsôṛô)}{বড়সড় (bôṛôsôṛô)}]{\bn{বড়সড়} (bôṛôsôṛô) — fairly big}

\label{lesson:BN-C143-borosoro}



\begin{quote}
Before the new one: say the Bengali for powerful, then the Bengali for tasty.
\end{quote}

\subsection*{You'll want to know: \bn{বড়সড়}}
\textbf{\bn{বড়সড়}} — \emph{bôṛôsôṛô} — \textquotedblleft{}fairly big\textquotedblright{}.

\begin{grammarlens}[title={What you've built}]
One more word for this chapter.
\end{grammarlens}

\subsection*{Your turn}
\begin{itemize}
\raggedright
  \item \emph{Say it:} \emph{bôṛôsôṛô}
  \item \emph{Say it:} \emph{bôṛôsôṛô}, once more
  \item \emph{Say it:} say \emph{bôṛôsôṛô}
  \item \emph{Recall:} say the Bengali for powerful, then the Bengali for tasty, then say \emph{bôṛôsôṛô} again
\end{itemize}

\begin{tcolorbox}[breakable,colback=teal!4,colframe=teal!35!black,title={Before you move on}]
What does \bn{বড়সড়} mean? (\textquotedblleft{}Fairly big\textquotedblright{}.) Say it once more.
\end{tcolorbox}
\section[\texorpdfstring{\bn{ছোটখাটো} (chhoṭôkhāṭo)}{ছোটখাটো (chhoṭôkhāṭo)}]{\bn{ছোটখাটো} (chhoṭôkhāṭo) — small, minor}

\label{lesson:BN-C143-chhotokhato}



\begin{quote}
Before the new one: say the Bengali for tasty, then the Bengali for fairly big.
\end{quote}

\subsection*{You'll want to know: \bn{ছোটখাটো}}
\textbf{\bn{ছোটখাটো}} — \emph{chhoṭôkhāṭo} — \textquotedblleft{}small, minor\textquotedblright{}.

\begin{grammarlens}[title={What you've built}]
One more word for this chapter.
\end{grammarlens}

\subsection*{Your turn}
\begin{itemize}
\raggedright
  \item \emph{Say it:} \emph{chhoṭôkhāṭo}
  \item \emph{Say it:} \emph{chhoṭôkhāṭo}, once more
  \item \emph{Say it:} say \emph{chhoṭôkhāṭo}
  \item \emph{Recall:} say the Bengali for tasty, then the Bengali for fairly big, then say \emph{chhoṭôkhāṭo} again
\end{itemize}

\begin{tcolorbox}[breakable,colback=teal!4,colframe=teal!35!black,title={Before you move on}]
What does \bn{ছোটখাটো} mean? (\textquotedblleft{}Small, minor\textquotedblright{}.) Say it once more.
\end{tcolorbox}
\section[\texorpdfstring{\bn{হাসিখুশি} (hāsikhushi)}{হাসিখুশি (hāsikhushi)}]{\bn{হাসিখুশি} (hāsikhushi) — cheerful}

\label{lesson:BN-C143-hasikhushi}



\begin{quote}
Before the new one: say the Bengali for fairly big, then the Bengali for small.
\end{quote}

\subsection*{You'll want to know: \bn{হাসিখুশি}}
\textbf{\bn{হাসিখুশি}} — \emph{hāsikhushi} — \textquotedblleft{}cheerful\textquotedblright{}.

\begin{grammarlens}[title={What you've built}]
One more word for this chapter.
\end{grammarlens}

\subsection*{Your turn}
\begin{itemize}
\raggedright
  \item \emph{Say it:} \emph{hāsikhushi}
  \item \emph{Say it:} \emph{hāsikhushi}, once more
  \item \emph{Say it:} say \emph{hāsikhushi}
  \item \emph{Recall:} say the Bengali for fairly big, then the Bengali for small, then say \emph{hāsikhushi} again
\end{itemize}

\begin{tcolorbox}[breakable,colback=teal!4,colframe=teal!35!black,title={Before you move on}]
What does \bn{হাসিখুশি} mean? (\textquotedblleft{}Cheerful\textquotedblright{}.) Say it once more.
\end{tcolorbox}
